\textbf{结论名称：}抽象函数定义域：定义域始终是 \(x\) 的范围。  

\textbf{适用条件：}已知 \(f(x)\) 或 \(f(g(x))\) 的定义域，求与同一对应法则相关的另一函数定义域。  

\textbf{变量说明：} \(x\)——自变量； \(f\)——对应关系； \(g(x)\)——内层函数表达式； \(D\)、\(A\) 分别表示已知定义域集合。  

\textbf{核心公式：}  
\[
\boxed{
\begin{aligned}
&\text{若 } f(x) \text{ 的定义域为 } D,\text{ 则 } f[g(x)] \text{ 的定义域为 } \{x \mid g(x) \in D\};\\
&\text{若 } f[g(x)] \text{ 的定义域为 } A,\text{ 则 } f(x) \text{ 的定义域为 } \{g(x) \mid x \in A\}.
\end{aligned}
}
\]  

\textbf{使用场景：}快速进行抽象函数定义域的互推，避免混淆中间变量与自变量的取值范围。  

\textbf{简单例子：}  
(1) 已知 \(f(x)\) 定义域为 \([0,2]\)，求 \(f(2x-1)\) 的定义域。  
由 \(0 \le 2x-1 \le 2\) 解得 \(\frac{1}{2} \le x \le \frac{3}{2}\)，定义域为 \(\big[\frac12,\frac32\big]\)。  
(2) 已知 \(f(2x-1)\) 定义域为 \([0,1]\)，求 \(f(x)\) 的定义域。  
当 \(x\in[0,1]\) 时，\(2x-1\in[-1,1]\)，故 \(f(x)\) 定义域为 \([-1,1]\)。